\section*{S4 Exact solvers for the subsample problem}
\label{supp:solvers}

We compared four exact solvers on the same subsamples, $30$ for each design and size, drawn from
simulated data sets with $10\%$ of outliers (Section~4.1 of the paper), each solver on one core
and all on the problem that Algorithm~\ref{alg:adaptive} needs: \eqref{eq:lsfit} with every
group of at least $d+1$ units. The first is the depth-first search of Section~S7, which also prunes a
partial labelling when the units left cannot bring every group to $d+1$ units, without and with
an incumbent from ten starts of hard alternation. The second is repetitive branch and bound: the
unconstrained problems on the last $1,2,\dots,m-1$ units are solved in turn, and the optimal value
of the units not yet labelled is added to the bound, which is valid because the residual sum of
squares of a group is superadditive over disjoint sets of its units. The third is the mixed
logical-quadratic programme of \citet{carbonneau2011globally}, with binary labels $z_{ik}$, free
coefficients, residuals $e_i$ with $z_{ik}=1\Rightarrow e_i=y_i-x_i^{\top}\theta_k$ as indicator
constraints, the quadratic objective $\sum_ie_i^{2}$, group sizes $\sum_iz_{ik}\ge d+1$ and the
symmetry-breaking constraints $z_{ik}=0$ for $k>i$, solved by Gurobi~13.0 \citep{gurobi2026}.
The fourth is the same programme with big-$M$ constraints, which need bounded coefficients
($|\theta_{kj}|\le50$, $M_i=|y_i|+50\sum_j|x_{ij}|$), and the objective written as an epigraph
constraint, solved by the open-source SCIP~10.0 through PySCIPOpt \citep{maher2016pyscipopt}.
The limit was $60$ seconds for the programmes and $2\times10^{8}$ nodes, about a minute, for the
searches. The probabilistic branch and bound of \citet{fois2026clusterwise} targets whole samples
without outliers, and we did not run it.

\begin{table}[htbp]
\centering
\small
\caption{Exact solution of the admissible subsample problem: median time in milliseconds over
$30$ subsamples, and the median number of nodes of our search. A superscript counts the solves
that reached the limit.}
\label{tab:solvers}
\begin{tabular}{lccrrrrrr}
\toprule
 & & & \multicolumn{3}{c}{Branch and bound} & \multicolumn{2}{c}{MIQP} & \\
\cmidrule(lr){4-6}\cmidrule(lr){7-8}
Design & $K$ & $m$ & ours & with incumbent & repetitive & Gurobi & SCIP & Nodes, ours \\
\midrule
D1 & 2 & 8 & 0.09 & 1.56 & 0.23 & 7.60 & 426 & 87 \\
D1 & 2 & 12 & 0.20 & 1.63 & 0.70 & 17.9 & 698 & 356 \\
D1 & 2 & 16 & 0.39 & 1.59 & 1.49 & 29.3 & 942 & 916 \\
D1 & 2 & 20 & 1.26 & 1.84 & 4.57 & 44.5 & 1\,617 & 4\,142 \\
D1 & 2 & 24 & 2.80 & 2.65 & 10.2 & 72.0 & 2\,846 & 10\,179 \\
D1 & 2 & 28 & 8.94 & 4.66 & 21.6 & 140 & 3\,783 & 33\,713 \\
D1 & 2 & 32 & 70.5$^{1}$ & 35.5$^{1}$ & 171$^{1}$ & 176 & 5\,353 & 273\,510 \\
D1 & 2 & 40 & 736$^{5}$ & 146$^{4}$ & 1\,473$^{4}$ & 317 & 10\,496 & 3\,168\,871 \\
D2 & 3 & 12 & 0.84 & 2.53 & 2.24 & 95.0 & 2\,670 & 1\,952 \\
D2 & 3 & 16 & 3.46 & 3.83 & 10.1 & 320 & 9\,774 & 10\,761 \\
D2 & 3 & 20 & 21.0 & 20.9 & 44.9 & 666 & 23\,163 & 75\,720 \\
D2 & 3 & 24 & 66.0 & 54.1 & 171 & 736 & 21\,211$^{1}$ & 346\,562 \\
D7 & 2 & 12 & 0.69 & 1.56 & 2.03 & 29.9 & 478 & 621 \\
D7 & 2 & 16 & 1.43 & 2.14 & 6.00 & 77.4 & 1\,590 & 2\,052 \\
four lines & 4 & 16 & 11.2 & 9.46 & 27.7 & 1\,507 & 13\,855 & 47\,284 \\
four lines & 4 & 20 & 37.2 & 25.1 & 120 & 6\,527 & 35\,404$^{2}$ & 166\,304 \\
\bottomrule
\end{tabular}
\end{table}

Whenever two solvers finished, their optimal values agreed to a relative $2\times10^{-6}$, except
that SCIP returned a larger value on $42$ of the $477$ subsamples it solved: there the optimum has a group fitted by a line
with a coefficient beyond $\pm50$, which the bound of the big-$M$ programme excludes.
Table~\ref{tab:solvers} gives the times. At the sizes that ESF uses, $m\le20$, the
depth-first search was $30$ to $180$ times faster than Gurobi and $700$ to $4{,}700$ times faster
than SCIP; the SCIP programme differs in its big-$M$ form and epigraph objective, so this factor
reflects the formulation as well as the solver. The advantage shrinks as $m$ grows, because the search is exponential in $m$ while the
programme's relaxations prune better on larger problems: with two groups, Gurobi was $16$ times
slower at $m=28$, $2.5$ times slower at $m=32$, and faster from $m=40$, where the search reached
its node limit on $5$ of $30$ subsamples. An incumbent from alternation pays from $m=24$ with two groups and from $m=16$ with four, and the
repetitive bound made the search two to four times slower at every size: solving the nested
subproblems costs more than the bound prunes. With about $110$ solves per data set, a mixed-integer
solver would turn the milliseconds of Algorithm~\ref{alg:adaptive} into seconds or minutes. Table~\ref{tab:S26} extends the comparison. With four covariates the
search was $30$ to $33$ times faster than Gurobi; with five groups and $m=20$ it took a median of
$0.19$ seconds, while Gurobi reached its time limit on $25$ of $30$ subsamples; and starting each
nested problem of the repetitive scheme from the optimum of the next smaller one, as
\citet{brusco2006repetitive} does, made it faster at the larger sizes but never as fast as the
plain search. The
subsample size of Algorithm~\ref{alg:adaptive} thus lies where the simple search is the right
tool, and a larger $m$, which Section~2.5 of the paper shows does not help, would call for a
mixed-integer solver.
